%=====================================================================
\chapter{The Semester Project}\label{ch:project}
%=====================================================================
\locator{6}{Part VI --- Reach, Law and Practice}

\begin{outcomes}
By the end of this chapter you will be able to:
\begin{tight}
  \item \textbf{Assemble} the shop you have built since \chref{stack} into a
        deliverable with a stated scope.
  \item \textbf{Deploy} it from nothing, on a machine that is not yours, in one
        command.
  \item \textbf{Measure} your own shop and \textbf{report} the figures
        honestly, including the ones that did not improve.
  \item \textbf{Present} and \textbf{defend} the decisions you made, naming
        what you would do differently.
  \item \textbf{Judge} another team's shop against the same criteria.
\end{tight}
\end{outcomes}

\begin{prereq}
Everything. Specifically: a shop that stands up (\chref{stack}), takes an order
(\chref{checkout}), charges correctly (\chref{shipping}), is hardened
(\chref{security}), carries at least one module you wrote (\chref{modules}),
and can be restored from backup (\chref{deploy}).
\end{prereq}

\begin{keyterms}
scope $\bullet$ deliverable $\bullet$ reproducible build $\bullet$ baseline
$\bullet$ evidence $\bullet$ defence $\bullet$ peer review $\bullet$ handover
$\bullet$ known limitation
\end{keyterms}

\section{What Is Being Assessed}

HEC's own teaching methodology for this course names \emph{Project} and
\emph{Presentations}, and this chapter discharges both. The project is not a
new piece of work: \textbf{it is the shop you have been adding to since week
three}, finished, deployed, measured and defended.

\begin{conceptbox}[The one sentence the whole project is marked against]
\textbf{Can somebody else bring your shop up from nothing, place an order in
it, and check your numbers?}

\smallskip
Everything below follows from that. A shop that only runs on your laptop is not
deployed. A claim you cannot reproduce is not a result. A decision you cannot
explain is not a decision.
\end{conceptbox}

\section{The Deliverable}

\begin{examplebox}[Five things, and nothing else]
\begin{tightnum}
  \item \textbf{A repository that builds.} \texttt{docker compose up -d
        -{}-build} on a clean machine produces a working shop with your
        catalogue in it. If a human step is needed, it is one line in the
        README and it is tested.
  \item \textbf{The shop itself}: a real catalogue of at least thirty products
        with options, categories, prices, weights and images; working shipping
        and tax for a stated place; a completable checkout; at least one
        payment method; and the \chref{security} hardening applied.
  \item \textbf{At least one module you wrote} (\chref{modules}), with its
        admin half, its language file, and an install step --- and a sentence
        saying what the shop could not do without it.
  \item \textbf{A measurement report}, two to four pages: what you measured,
        with what, on what machine, what you found, and what you changed as a
        result. Figures you did not improve are part of the report.
  \item \textbf{A ten-minute presentation and a defence.}
\end{tightnum}

\smallskip
\textbf{Not wanted:} a redesigned theme as the main contribution, a slide deck
longer than the shop, or a feature list. The marks are in whether it works,
whether it is measured, and whether you can defend it.
\end{examplebox}

\rowcolors{2}{white}{lightbg}
\begin{longtable}{@{}L{5.6cm}L{2.0cm}L{6.6cm}@{}}
\toprule
\rowcolor{primary!15}
\textbf{Assessed on} & \textbf{Weight} & \textbf{What earns it} \\
\midrule
\endhead
It builds and runs elsewhere & 20\% & Somebody else runs one command \\
The shop works end to end & 20\% & An order can be placed and fulfilled \\
Your own extension & 15\% & It installs, removes cleanly, and is needed \\
Hardening and correctness & 15\% & The \chref{security} battery, and the
go-live list from \chref{deploy} \\
Measurement report & 20\% & Reproducible figures, honestly reported \\
Presentation and defence & 10\% & You can answer ``why'' \\
\bottomrule
\end{longtable}

\section{The Measurement Report}

This is where the course differs from a web-development project, and it is
worth being explicit about what a good one looks like.

\begin{conceptbox}[Four measurements, chosen because you have already made them
once]
\begin{tightnum}
  \item \textbf{A catalogue measurement} (\chref{catalogue}): a page's time and
        query count against catalogue size, with and without an index you
        chose. \emph{Report the case where the index did not help, if there is
        one.}
  \item \textbf{A page weight measurement} (\chref{ux}): bytes and requests,
        before and after one change, with the change named.
  \item \textbf{A security measurement} (\chref{security}): the battery before
        and after your hardening, with the findings you could not close listed
        and costed.
  \item \textbf{One of your own choosing}, with a stated question it answers.
\end{tightnum}

\smallskip
\textbf{Every figure must name the file that produces it}, exactly as this
handout does. A number with no command behind it is an assertion.
\end{conceptbox}

\begin{alertbox}[What a good report does that a weak one does not]
\textbf{It states the machine.} Appendix~A of this handout exists because every
timing in it is a fact about one laptop. Yours is a fact about yours.

\smallskip
\textbf{It reports the measurement that disagreed with you.} This handout
printed three wrong versions of \chref{catalogue}'s experiment, the store
identifier that made 50{,}000 products invisible, and three security controls
that were applied correctly and did nothing. Every one of those is worth more
than a result that went as expected, because it shows the method working.

\smallskip
\textbf{It separates what it measured from what it believes.} ``The index made
the page slower at 50{,}000 products on this machine'' is a measurement. ``Our
shop is fast'' is not.

\smallskip
\textbf{It says what it could not measure.} \chref{ux}'s harness could not see
fonts requested by CSS, and said so. Knowing the limits of your instrument is
the difference between a report and a brochure.
\end{alertbox}

\section{The Defence}

Ten minutes of presentation and then questions. The questions are not a test of
memory.

\begin{examplebox}[What you will be asked, and what a good answer sounds like]
\begin{tight}
  \item \textbf{``Why this platform?''} Not ``it was in the course''.
        \chref{planning} is the argument; make it yours, with the alternative
        you rejected and why.
  \item \textbf{``Show me an order going through.''} Live, on the running shop.
        Know what status it lands in and what moved in the database when it
        did (\chref{checkout}, \chref{orders}).
  \item \textbf{``What does this number mean?''} Pointing at any figure in your
        report. Be able to say what it counts, what it excludes, and what would
        change it.
  \item \textbf{``What is still wrong with it?''} The best question and the one
        most teams are unprepared for. A team that names three real
        limitations is further along than one that claims none.
  \item \textbf{``What would you do with another month?''} Ordered by measured
        benefit, not by what would be interesting to build.
  \item \textbf{``Who did what?''} Answered the same way by every member.
\end{tight}
\end{examplebox}

\begin{conceptbox}[Reviewing another team's shop]
Each team reviews one other, using the same instruments, and submits one page.

\begin{tightnum}
  \item \textbf{Clone and build it.} Time how long until you see a storefront.
        That number is the review's first finding.
  \item \textbf{Place an order.} Note every step that confused you.
  \item \textbf{Run the security battery} against \emph{their} container on
        \emph{your} machine. This is within scope precisely because you have
        built it yourself and it is running on your own computer ---
        \chref{security}'s rule is about systems you do not control, and a
        classmate's repository running on your laptop is one you do.
  \item \textbf{Check one of their numbers.} Run the file they name, and report
        what you got.
  \item \textbf{Write one page}: what works, what you could not get working,
        one thing you will steal, and one thing you would change.
\end{tightnum}

\smallskip
\textbf{Reproducing somebody else's number is the single most valuable hour of
the semester}, in both directions. It is also the only part of this course that
cannot be done alone.
\end{conceptbox}

\section{A Realistic Schedule}

\rowcolors{2}{white}{lightbg}
\begin{longtable}{@{}L{2.4cm}L{5.4cm}L{6.4cm}@{}}
\toprule
\rowcolor{primary!15}
\textbf{Week} & \textbf{Milestone} & \textbf{What is actually due} \\
\midrule
\endhead
10 & Proposal & One page: what you sell, to whom, where you deliver, which
payment methods, and the one feature you will build \\
12 & Catalogue complete & 30+ real products with weights and images, and your
first measurement \\
14 & Your module works & Installs and removes cleanly on a clean container \\
15 & Hardened and deployed & The go-live list from \chref{deploy}, and a
restore you have timed \\
16 & Present and defend & And review another team \\
\bottomrule
\end{longtable}

\begin{pitfall}
\begin{tight}
  \item \textbf{Leaving deployment to the last week.} It is the deliverable
        with the most ways to fail and the least room to recover.
  \item \textbf{A shop that runs only on the machine it was built on.}
  \item \textbf{Measuring once, at the end.} A baseline you did not take in
        week 10 cannot be taken in week 16.
  \item \textbf{Reporting only the figures that improved.}
  \item \textbf{Quoting a number without the command that produces it.}
  \item \textbf{Spending the project on the theme.} It is the most visible work
        and the least of the marks.
  \item \textbf{Twelve demo products with no weights}, which breaks shipping
        in a way \chref{shipping} told you about.
  \item \textbf{A module that cannot be removed}, which means it was never
        tested from a clean state.
  \item \textbf{Arriving at the defence unable to say what is still wrong with
        it.}
  \item \textbf{Running anybody's security battery against anything but your
        own or your classmate's container on your own machine}
        (\chref{security}).
\end{tight}
\end{pitfall}

\begin{lab}[Ship it]
\begin{lstlisting}[language=bash]
# --- 1. Prove it builds from nothing --------------------------------
cd code
docker compose down -v
docker compose up -d --build
python waitfor.py
# Time this. It is the first number in your report.

# --- 2. Place an order on the fresh shop ----------------------------
python m10_checkout.py --repeat 1
python m10_checkout.py --states
# If this fails on a clean build, your README is wrong.

# --- 3. Install your own extension ----------------------------------
python m19_modules.py --install
python m19_modules.py --measure

# --- 4. Harden, and record the battery before and after -------------
python m16_security.py

# --- 5. Take the four measurements ----------------------------------
python m06_catalogue.py --quick
python m08_pageweight.py --compare
python m20_deploy.py --load
# Plus your own. Save every output; it is your evidence.

# --- 6. Back it up, destroy it, restore it --------------------------
python m20_deploy.py --backup
# Then do it for real: down -v, up, restore, and check the catalogue
# is there. Record the restore time.

# --- 7. Write the README somebody else will follow ------------------
# Hand it to a classmate and watch them. Every question they ask is
# a line missing from it. Do not help them; write it down.

# --- 8. Write the report --------------------------------------------
# Four measurements, the machine, what changed, what did not, and
# what you could not measure.

# --- 9. Prepare three limitations -----------------------------------
# Real ones, with what each would cost to fix. You will be asked.
\end{lstlisting}

\smallskip
\textbf{What to notice.} Step 7 is the one that decides the first 20\% of the
marks, and watching somebody fail to follow your README is the only way to
write one. Every course does this badly; this is your chance not to.
\end{lab}

\begin{checkpoint}
\begin{tightnum}
  \item Your shop runs perfectly on your laptop and a classmate cannot get it
        started. List the five most likely causes in order, say how you would
        find out which it is without access to their machine, and name the one
        thing you should have done in week 12 that would have prevented all
        five.
  \item You measured your category page at 180\,ms before an index and 195\,ms
        after. Say what you would put in the report, what you would investigate
        first, and why this result is worth more marks than one that went the
        way you expected.
  \item At the defence you are asked ``what is still wrong with it?''. Give
        three answers for the shop as this handout left it, each with the
        chapter that found it and the cost of fixing it.
\end{tightnum}
\end{checkpoint}

\begin{chaptersummary}
\begin{tight}
  \item The project is the shop you have built since \chref{stack}, finished
        and defended --- not a new piece of work.
  \item One sentence decides it: can somebody else bring it up from nothing,
        place an order, and check your numbers?
  \item Five deliverables: a repository that builds, a shop that works end to
        end, an extension you wrote, a measurement report, and a defence.
  \item Every figure names the file that produces it, every timing names the
        machine, and the measurements that disagreed with you are worth more
        than the ones that did not.
  \item A good report separates what it measured from what it believes, and
        says what its instrument could not see.
  \item The hardest defence question is ``what is still wrong with it?''. A
        team naming three real limitations is ahead of one claiming none.
  \item Reproducing another team's number is the most valuable hour of the
        semester, and the only part of the course that cannot be done alone.
  \item Deploy early. A baseline not taken in week 10 cannot be taken in week
        16.
\end{tight}
\end{chaptersummary}

\begin{reviewq}
  \item \textbf{(MCQ)} The largest single share of the project mark is for:
        (a)~the theme; (b)~the number of features; (c)~that it builds and runs
        on another machine; (d)~the length of the presentation.
  \item \textbf{(MCQ)} A measurement in your report that did not improve
        should be: (a)~omitted; (b)~re-run until it does; (c)~reported with
        what you think caused it; (d)~replaced with an estimate.
  \item \textbf{(Short)} Explain why a baseline must be taken early, with an
        example from this course.
  \item \textbf{(Short)} Give two things a measurement report must state about
        the machine, and say why each matters.
  \item \textbf{(Short)} Why is being able to name your shop's limitations
        worth marks, rather than costing them?
  \item \textbf{(Applied)} Write the week-10 proposal for a shop selling a
        product of your choosing in Pakistan: what you sell, to whom, where you
        deliver, which payment methods and why, the tax configuration, and the
        one feature you will build with the reason the shop needs it.
  \item \textbf{(Applied)} You have one month after the semester to take the
        shop to real customers. Give your ordered plan, each step justified by
        a measurement or a chapter from this course, and state the three things
        you would refuse to launch without.
\end{reviewq}
